% GENERATED FILE. Edit canonical lessons, then run npm run generate:books.
% canonical-source-hash: fnv1a64:5435393faced533a
% canonical-lessons: KA-C257-nalli, KA-C257-mettilu, KA-C257-hittu, KA-C257-bele, KA-C257-hunasehannu

\chapter{Tap, Stairs, Flour, Lentils, Tamarind}
\label{ch:ka-tap-stairs-flour-lentils-tamarind}
\hlchaptermodality{257}

\begin{chapteropening}
\textbf{By the end of this chapter:} \emph{I can say a tap, stairs, flour, lentils and tamarind.}

Complete the last lesson of chapter 257: I can say a tap, stairs, flour, lentils and tamarind.
\end{chapteropening}

\section[\texorpdfstring{\kn{ನಲ್ಲಿ} (nalli)}{ನಲ್ಲಿ (nalli)}]{\kn{ನಲ್ಲಿ} (nalli) — a tap}

\label{lesson:KA-C257-nalli}



\begin{quote}
Before the new one: say the Kannada for blurred, then the Kannada for crowded.
\end{quote}

\subsection*{You'll want to know: \kn{ನಲ್ಲಿ}}
\textbf{\kn{ನಲ್ಲಿ}} — \emph{nalli} — \textquotedblleft{}a tap\textquotedblright{}.

\begin{grammarlens}[title={What you've built}]
One more word for this chapter.
\end{grammarlens}

\subsection*{Your turn}
\begin{itemize}
\raggedright
  \item \emph{Say it:} \emph{nalli}
  \item \emph{Say it:} \emph{nalli}, once more
  \item \emph{Say it:} say \emph{nalli}
  \item \emph{Recall:} say the Kannada for blurred, then the Kannada for crowded, then say \emph{nalli} again
\end{itemize}

\begin{tcolorbox}[breakable,colback=teal!4,colframe=teal!35!black,title={Before you move on}]
What does \kn{ನಲ್ಲಿ} mean? (\textquotedblleft{}A tap\textquotedblright{}.) Say it once more.
\end{tcolorbox}
\section[\texorpdfstring{\kn{ಮೆಟ್ಟಿಲು} (meṭṭilu)}{ಮೆಟ್ಟಿಲು (meṭṭilu)}]{\kn{ಮೆಟ್ಟಿಲು} (meṭṭilu) — stairs, a step}

\label{lesson:KA-C257-mettilu}



\begin{quote}
Before the new one: say the Kannada for crowded, then the Kannada for a tap.
\end{quote}

\subsection*{You'll want to know: \kn{ಮೆಟ್ಟಿಲು}}
\textbf{\kn{ಮೆಟ್ಟಿಲು}} — \emph{meṭṭilu} — \textquotedblleft{}stairs, a step\textquotedblright{}.

\begin{grammarlens}[title={What you've built}]
One more word for this chapter.
\end{grammarlens}

\subsection*{Your turn}
\begin{itemize}
\raggedright
  \item \emph{Say it:} \emph{meṭṭilu}
  \item \emph{Say it:} \emph{meṭṭilu}, once more
  \item \emph{Say it:} say \emph{meṭṭilu}
  \item \emph{Recall:} say the Kannada for crowded, then the Kannada for a tap, then say \emph{meṭṭilu} again
\end{itemize}

\begin{tcolorbox}[breakable,colback=teal!4,colframe=teal!35!black,title={Before you move on}]
What does \kn{ಮೆಟ್ಟಿಲು} mean? (\textquotedblleft{}Stairs, a step\textquotedblright{}.) Say it once more.
\end{tcolorbox}
\section[\texorpdfstring{\kn{ಹಿಟ್ಟು} (hiṭṭu)}{ಹಿಟ್ಟು (hiṭṭu)}]{\kn{ಹಿಟ್ಟು} (hiṭṭu) — flour}

\label{lesson:KA-C257-hittu}



\begin{quote}
Before the new one: say the Kannada for a tap, then the Kannada for stairs.
\end{quote}

\subsection*{You'll want to know: \kn{ಹಿಟ್ಟು}}
\textbf{\kn{ಹಿಟ್ಟು}} — \emph{hiṭṭu} — \textquotedblleft{}flour\textquotedblright{}.

\begin{grammarlens}[title={What you've built}]
One more word for this chapter.
\end{grammarlens}

\subsection*{Your turn}
\begin{itemize}
\raggedright
  \item \emph{Say it:} \emph{hiṭṭu}
  \item \emph{Say it:} \emph{hiṭṭu}, once more
  \item \emph{Say it:} say \emph{hiṭṭu}
  \item \emph{Recall:} say the Kannada for a tap, then the Kannada for stairs, then say \emph{hiṭṭu} again
\end{itemize}

\begin{tcolorbox}[breakable,colback=teal!4,colframe=teal!35!black,title={Before you move on}]
What does \kn{ಹಿಟ್ಟು} mean? (\textquotedblleft{}Flour\textquotedblright{}.) Say it once more.
\end{tcolorbox}
\section[\texorpdfstring{\kn{ಬೇಳೆ} (bēḷe)}{ಬೇಳೆ (bēḷe)}]{\kn{ಬೇಳೆ} (bēḷe) — lentils, dal}

\label{lesson:KA-C257-bele}



\begin{quote}
Before the new one: say the Kannada for stairs, then the Kannada for flour.
\end{quote}

\subsection*{You'll want to know: \kn{ಬೇಳೆ}}
\textbf{\kn{ಬೇಳೆ}} — \emph{bēḷe} — \textquotedblleft{}lentils, dal\textquotedblright{}.

\begin{grammarlens}[title={What you've built}]
One more word for this chapter.
\end{grammarlens}

\subsection*{Your turn}
\begin{itemize}
\raggedright
  \item \emph{Say it:} \emph{bēḷe}
  \item \emph{Say it:} \emph{bēḷe}, once more
  \item \emph{Say it:} say \emph{bēḷe}
  \item \emph{Recall:} say the Kannada for stairs, then the Kannada for flour, then say \emph{bēḷe} again
\end{itemize}

\begin{tcolorbox}[breakable,colback=teal!4,colframe=teal!35!black,title={Before you move on}]
What does \kn{ಬೇಳೆ} mean? (\textquotedblleft{}Lentils, dal\textquotedblright{}.) Say it once more.
\end{tcolorbox}
\section[\texorpdfstring{\kn{ಹುಣಸೆಹಣ್ಣು} (huṇasehaṇṇu)}{ಹುಣಸೆಹಣ್ಣು (huṇasehaṇṇu)}]{\kn{ಹುಣಸೆಹಣ್ಣು} (huṇasehaṇṇu) — tamarind}

\label{lesson:KA-C257-hunasehannu}



\begin{quote}
Before the new one: say the Kannada for flour, then the Kannada for lentils.
\end{quote}

\subsection*{You'll want to know: \kn{ಹುಣಸೆಹಣ್ಣು}}
\textbf{\kn{ಹುಣಸೆಹಣ್ಣು}} — \emph{huṇasehaṇṇu} — \textquotedblleft{}tamarind\textquotedblright{}.

\begin{grammarlens}[title={What you've built}]
One more word for this chapter.
\end{grammarlens}

\subsection*{Your turn}
\begin{itemize}
\raggedright
  \item \emph{Say it:} \emph{huṇasehaṇṇu}
  \item \emph{Say it:} \emph{huṇasehaṇṇu}, once more
  \item \emph{Say it:} say \emph{huṇasehaṇṇu}
  \item \emph{Recall:} say the Kannada for flour, then the Kannada for lentils, then say \emph{huṇasehaṇṇu} again
\end{itemize}

\begin{tcolorbox}[breakable,colback=teal!4,colframe=teal!35!black,title={Before you move on}]
What does \kn{ಹುಣಸೆಹಣ್ಣು} mean? (\textquotedblleft{}Tamarind\textquotedblright{}.) Say it once more.
\end{tcolorbox}
